
\documentclass[12pt,openany]{book}
\usepackage{amsmath}
\usepackage{amssymb}
\usepackage{amsthm}
\usepackage{mathtools}
\usepackage[hidelinks]{hyperref}
\usepackage{booktabs}
\usepackage{geometry}
\usepackage{fancyhdr}
\usepackage{setspace}
\usepackage[expansion=false,nopatch=footnote]{microtype}
\usepackage{xcolor}
\usepackage{enumitem}
\usepackage{titlesec}
\usepackage{tocloft}
\usepackage{array}
\usepackage{tabularx}
\usepackage{float}
\usepackage{longtable}
\usepackage{ragged2e}
\usepackage{thmtools}
\usepackage{chngcntr}
\counterwithout{table}{chapter}
\counterwithout{figure}{chapter}
\usepackage{sectionbreak}
\usepackage{algorithmicx}
\usepackage{algpseudocode}
\usepackage{algorithm}
\usepackage{url}
\Urlmuskip=0mu plus 2mu\relax

% ============================================================
% CUSTOM MACROS
% ============================================================
\newcommand{\R}{\mathbb{R}}
\newcommand{\E}{\mathbb{E}}
\newcommand{\Prob}{\mathbb{P}}
\newcommand{\N}{\mathbb{N}}

\newcommand{\bfw}{\mathbf{w}}
\newcommand{\bfx}{\mathbf{x}}
\newcommand{\bfz}{\mathbf{z}}
\newcommand{\bfy}{\mathbf{y}}
\newcommand{\bfmu}{\boldsymbol{\mu}}
\newcommand{\bfSigma}{\boldsymbol{\Sigma}}
\newcommand{\bfQ}{\mathbf{Q}}
\newcommand{\bfR}{\mathbf{R}}
\newcommand{\bfF}{\mathbf{F}}
\newcommand{\bfH}{\mathbf{H}}
\newcommand{\bfK}{\mathbf{K}}
\newcommand{\bfP}{\mathbf{P}}
\newcommand{\bfI}{\mathbf{I}}
\newcommand{\bfA}{\mathbf{A}}
\newcommand{\bfe}{\mathbf{e}}
\newcommand{\bfv}{\mathbf{v}}
\newcommand{\bfu}{\mathbf{u}}
\newcommand{\bfGamma}{\boldsymbol{\Gamma}}
\newcommand{\bfOmega}{\boldsymbol{\Omega}}
\newcommand{\bfxi}{\boldsymbol{\xi}}

\newcommand{\calE}{\mathcal{E}}
\newcommand{\calF}{\mathcal{F}}
\newcommand{\calG}{\mathcal{G}}
\newcommand{\calP}{\mathcal{P}}
\newcommand{\calU}{\mathcal{U}}
\newcommand{\calM}{\mathcal{M}}
\newcommand{\calT}{\mathcal{T}}
\newcommand{\calR}{\mathcal{R}}
\newcommand{\calS}{\mathcal{S}}
\newcommand{\calI}{\mathcal{I}}
\newcommand{\calX}{\mathbf{X}}

\newcommand{\MaxDD}{\mathrm{MaxDD}}
\newcommand{\DD}{\mathrm{DD}}

\DeclareMathOperator{\CVaR}{CVaR}
\DeclareMathOperator{\VaR}{VaR}
\DeclareMathOperator{\Var}{Var}
\DeclareMathOperator{\Cov}{Cov}
\DeclareMathOperator{\tr}{tr}
\DeclareMathOperator{\rank}{rank}
\DeclareMathOperator*{\argmax}{arg\,max}
\DeclareMathOperator*{\argmin}{arg\,min}

% ============================================================
% GEOMETRY AND LAYOUT SETTINGS
% ============================================================
\geometry{margin=1in, includefoot}
\microtypesetup{expansion=false, final}
\singlespacing
\setlength{\parindent}{0.5in}
\setlength{\parskip}{0.5ex plus 0.2ex minus 0.1ex}
\setlength{\emergencystretch}{1.5em}
\tolerance=9999
\hbadness=10000
\vbadness=10000
\hyphenpenalty=10000
\exhyphenpenalty=10000
\raggedbottom
\sloppy

\setlength{\abovedisplayskip}{6pt plus 2pt minus 3pt}
\setlength{\belowdisplayskip}{6pt plus 2pt minus 3pt}
\setlength{\abovedisplayshortskip}{3pt plus 1pt minus 1pt}
\setlength{\belowdisplayshortskip}{3pt plus 1pt minus 1pt}
\numberwithin{equation}{chapter}
\allowdisplaybreaks[3]

% ============================================================
% THEOREM ENVIRONMENTS
% ============================================================
\theoremstyle{plain}
\newtheorem{theorem}{Theorem}[chapter]
\newtheorem{lemma}[theorem]{Lemma}
\newtheorem{corollary}[theorem]{Corollary}
\newtheorem{proposition}[theorem]{Proposition}
\newtheorem{conjecture}[theorem]{Conjecture}
\theoremstyle{definition}
\newtheorem{definition}[theorem]{Definition}
\newtheorem{axiom}[theorem]{Axiom}
\newtheorem{assumption}[theorem]{Assumption}
\theoremstyle{remark}
\newtheorem{remark}[theorem]{Remark}

\providecommand{\theHtheorem}{}
\renewcommand{\theHtheorem}{\thechapter-\arabic{theorem}}
\providecommand{\theHlemma}{}
\renewcommand{\theHlemma}{\thechapter-\arabic{theorem}}
\providecommand{\theHcorollary}{}
\renewcommand{\theHcorollary}{\thechapter-\arabic{theorem}}
\providecommand{\theHproposition}{}
\renewcommand{\theHproposition}{\thechapter-\arabic{theorem}}
\providecommand{\theHdefinition}{}
\renewcommand{\theHdefinition}{\thechapter-\arabic{theorem}}
\providecommand{\theHaxiom}{}
\renewcommand{\theHaxiom}{\thechapter-\arabic{theorem}}
\providecommand{\theHassumption}{}
\renewcommand{\theHassumption}{\thechapter-\arabic{theorem}}
\providecommand{\theHremark}{}
\renewcommand{\theHremark}{\thechapter-\arabic{theorem}}
\providecommand{\theHconjecture}{}
\renewcommand{\theHconjecture}{\thechapter-\arabic{theorem}}

\makeatletter
\providecommand*{\toclevel@conjecture}{0}
\makeatother

\newenvironment{abstract}{%
  \chapter*{Abstract}
  \begin{quote}\small
}{%
  \end{quote}
}
\newenvironment{keywords}{%
  \paragraph*{Keywords:}
}{}

% ============================================================
% SECTION FORMATTING
% ============================================================
\titleformat{\chapter}[display]
    {\normalfont\Huge\bfseries}{\chaptertitlename\ \thechapter}{20pt}{\Huge}
\titlespacing*{\chapter}{0pt}{50pt}{40pt}
\titleformat{\section}{\normalfont\Large\bfseries}{\thesection}{1em}{}
\titlespacing*{\section}{0pt}{3.5ex plus 1ex minus 0.2ex}{2.3ex plus 0.2ex}
\titleformat{\subsection}{\normalfont\large\bfseries}{\thesubsection}{1em}{}
\titlespacing*{\subsection}{0pt}{3.25ex plus 1ex minus 0.2ex}{1.5ex plus 0.2ex}
\titleformat{\subsubsection}{\normalfont\normalsize\bfseries}{\thesubsubsection}{1em}{}
\titlespacing*{\subsubsection}{0pt}{3ex plus 0.5ex minus 0.2ex}{1.5ex plus 0.2ex}

% ============================================================
% TABLE OF CONTENTS
% ============================================================
\renewcommand{\cftchapleader}{\cftdotfill{\cftdotsep}}
\setlength{\cftchapindent}{0em}
\setlength{\cftsecindent}{1.5em}

% ============================================================
% HEADER AND FOOTER
% ============================================================
\pagestyle{fancy}
\fancyhf{}
\fancyhead[LE]{\thepage}
\fancyhead[RO]{\thepage}
\fancyhead[RE]{\nouppercase{\leftmark}}
\fancyhead[LO]{\nouppercase{\rightmark}}
\renewcommand{\headrulewidth}{0.4pt}
\renewcommand{\footrulewidth}{0pt}
\setlength{\headheight}{27.2pt}

% ============================================================
% MISCELLANEOUS
% ============================================================
\setlist[enumerate]{leftmargin=*, partopsep=0.1ex}
\clubpenalty=10000
\widowpenalty=10000
\displaywidowpenalty=10000
\renewcommand{\arraystretch}{1.2}
\setlength{\tabcolsep}{4pt}

\pdfstringdefDisableCommands{%
  \def\mathcal#1{#1}%
}

% ============================================================
% TITLE
% ============================================================
\title{%
  Multi-Agent Quantitative Investment Architecture:\\
  A Unified Framework of Kalman Filtering,\\
  Seven-State Capital Gating, and\\
  Finite-Capital Optimal Allocation\\[0.6em]
  {\large Public Judges Whitepaper ---
  Alpaca Trading Hackathon 2026}%
}
\author{Hadrian Hu}
\date{\today \quad Version 2026.1.0.0}

% ============================================================
\begin{document}
% ============================================================

\maketitle

\tableofcontents
\listoftheorems[ignoreall,show={theorem,lemma,proposition,corollary,
  definition,axiom,assumption,conjecture}]

% ==============================================================
\begin{abstract}
This whitepaper presents the complete mathematical architecture
of a multi-agent, cross-market quantitative investment platform
developed for the Alpaca Trading Hackathon 2026.  The system
operates on a fixed initial capital of \$100{,}000 and
coordinates heterogeneous specialised agents across four
financial domains --- economic, financial, fiscal, and investment
--- to produce uncertainty-aware, survival-first capital
allocation decisions.

The intellectual core is a three-layer formal identification:

\begin{enumerate}[label=(\roman*), leftmargin=2.5em, itemsep=2pt,
                  parsep=0pt, topsep=3pt]
  \item \textbf{Kalman isomorphism.}  The capital allocation
        cycle is structurally isomorphic to a multi-rate Kalman
        filter bank.  The Kalman gain is the confidence-weighted
        capital deployment fraction; the posterior covariance is
        residual portfolio risk; the innovation signal is the
        aggregate agent signal deviation from prior forecast.
  \item \textbf{Seven-state gating.}  The hidden market state
        is a seven-dimensional State-of-Charge vector
        $\mathbf{X}_t = (\calE_t, \calF_t, \calG_t, \calP_t,
        \calU_t, \calM_t, \calT_t)^\top \in [0,1]^7$ covering
        the economic, financial, fiscal, portfolio, fundamental,
        market, and sector dimensions.  Capital deployment is
        gated by the product function $G(\mathbf{X}_t) =
        \prod_{d=1}^7 g_d(\mathcal{S}_t^{(d)})$, which forces
        zero deployment when any single dimension falls below
        its critical threshold.
  \item \textbf{MVUE optimality.}  Under the finite capital
        constraint, the optimal agentic trading policy is
        the minimum-variance unbiased estimator (MVUE) of the
        latent market state given partial, noisy, multi-domain
        observations.  The finite capital constraint is a
        precision instrument, not a handicap: it forces the
        system into the well-calibrated, survival-preserving
        operating regime.
\end{enumerate}

All results are stated as formal definitions, axioms,
assumptions, propositions, lemmas, and theorems with complete
proofs.  Every boundary, degenerate, and extremal case is
addressed explicitly.  The monograph is fully auditable by any
independent third-party reviewer.
\end{abstract}

\begin{keywords}
multi-agent systems; Kalman filter; state-space model;
prediction--correction cycle; Kalman gain; MVUE;
seven-state investment architecture; State-of-Charge analogy;
economic state; financial state; fiscal state; portfolio state;
fundamental state; market state; sector state; product gating;
circuit-breaker; capital deployment; finite capital constraint;
composite risk-adjusted objective; CVaR; maximum drawdown;
Bayesian agent reputation; ensemble signal aggregation;
disagreement inflation; IMM; Interacting Multiple Model;
Kelly criterion; recursive reinvestment; portfolio diversification;
survival-first priority; Alpaca hackathon 2026
\end{keywords}
% ==============================================================
\chapter*{Executive Summary}
\addcontentsline{toc}{chapter}{Executive Summary}
% ==============================================================

\noindent
The central research question addressed by this platform is:

\begin{quote}
\emph{Given all information legitimately available at time $t$,
what market regime is most probable, what risks currently
dominate, how uncertain is that assessment, and what portfolio
allocation, hedge, or trading action provides the best
risk-adjusted response under a finite initial capital of
\$100{,}000?}
\end{quote}

\noindent
The answer is formalised as the following dual statement:

\begin{align}
  \text{Optimal Agentic Trading}
    &\;\neq\; \max\!\bigl(\text{Raw Profit}\bigr),
  \label{eq:exec-neq} \\[4pt]
  \text{Optimal Agentic Trading}
    &\;=\; \text{MVUE of Latent Market State Given Partial,}
  \notag \\
    &\qquad \text{Noisy, Multi-Domain Observations.}
  \label{eq:exec-eq}
\end{align}

\noindent
\textbf{Priority ordering.}\enspace
The system respects the following strict lexicographic priority,
proved consistent in the companion proof documents:

\begin{equation}
  \text{Survival} \;\succ\; \text{Preservation}
  \;\succ\; \text{Intelligent Allocation}
  \;\succ\; \text{Controlled Risk}
  \;\succ\; \text{Compounding.}
  \label{eq:priority}
\end{equation}

\noindent
\textbf{Architecture overview.}\enspace
The platform integrates four tightly coupled formal layers:

\begin{enumerate}[label=\textbf{\arabic*.}, leftmargin=2em,
                  itemsep=3pt, parsep=0pt, topsep=4pt]
  \item \textbf{Multi-agent signal ensemble.}  $N \geq 2$
        specialised agents each emit an eight-channel output
        tuple $(s_i, c_i, u_i, d_i, p_i^+, p_i^-, \Delta t_i,
        r_i)$.  The confidence-weighted, uncertainty-discounted
        ensemble aggregate signal $S_t \in [-1,1]$ and
        disagreement index $D_t \geq 0$ serve as the primary
        observation vector for the Kalman correction step.
  \item \textbf{Kalman filter capital allocation.}  The
        investment Kalman gain $K_t = P_{t|t-1} /
        (P_{t|t-1} + R_t)$ maps agent confidence and prior
        portfolio risk into a capital deployment fraction, where
        $R_t = \sigma_{\mathrm{base}}^2 (1-\bar{c}_t)/\bar{c}_t
        + D_t^2 \sigma_{\mathrm{base}}^2$ is the
        confidence-and-disagreement-inflated measurement noise.
  \item \textbf{Seven-state gating architecture.}  The hidden
        state vector $\mathbf{X}_t \in [0,1]^7$ is estimated by
        seven specialised state agents.  The product gating
        function $G(\mathbf{X}_t) \in [0,1]$ enforces a hard
        circuit-breaker: $G = 0$ if and only if any single state
        falls below its critical threshold, halting all new
        capital deployment regardless of signal strength.
  \item \textbf{Bayesian reputation and recursive
        reinvestment.}  Agent weights $w_i(t)$ are updated via
        Beta-Bernoulli conjugate posteriors, converging
        geometrically to the true reliability vector.  Profits
        are recursively reinvested under the composite
        risk-adjusted objective $J = \alpha G - \beta D_{\mathrm{dd}}
        - \gamma V + \delta S_{\mathrm{surv}} + \varepsilon
        R_{\mathrm{risk}} - \cdots$, achieving monotone
        compounding in expectation.
\end{enumerate}

\noindent
\textbf{Key formal results.}\enspace
The table below summarises the principal theorems.

\begin{table}[H]
\centering
\caption{Principal results and their investment implications.}
\label{tab:exec-results}
\small
\begin{tabularx}{\linewidth}{@{} l X X @{}}
\toprule
\textbf{Result} & \textbf{Formal Guarantee}
  & \textbf{Investment Implication} \\
\midrule
Bijection $\phi$
  & $\phi : \mathfrak{K} \xrightarrow{\;\sim\;} \mathfrak{A}$
  & Kalman filter IS the optimal capital allocator. \\
Gain bounds
  & $K_t \in [0,1)$; strictly monotone in $R_t$, $P_{t|t-1}$
  & Deployment auto-scales with confidence and prior risk. \\
State invariant
  & $\mathcal{S}_{t+1}^{(d)} \in [0,1]$ a.s.\ for all $d$, $t$
  & State charges never escape the unit interval. \\
Zero-gate condition
  & $G = 0 \Leftrightarrow \exists\,d: \mathcal{S}_t^{(d)} \leq
    \mathcal{S}_{\min}^{(d)}$
  & One critical breach halts ALL new deployment. \\
Full-deployment condition
  & $G = 1 \Leftrightarrow \forall\,d: \mathcal{S}_t^{(d)} \geq
    \mathcal{S}_{\mathrm{full}}^{(d)}$
  & All seven dimensions must be healthy for full deployment. \\
Disagreement inflation
  & $D_t^2 \sigma_{\mathrm{base}}^2$ added to $R_t$
  & High disagreement shrinks Kalman gain; suppresses deployment. \\
Objective equivalence
  & $\argmax_\pi J(\pi) = \argmin_\pi \E[\|\mathbf{X}_t -
    \hat{\mathbf{X}}_t\|^2]$
  & Maximising $J$ IS minimising posterior MSE. \\
MVUE policy
  & $\pi^* = K_t \cdot C_{\mathrm{avail}} \cdot G(\mathbf{X}_t)$
    achieves Cram\'{e}r--Rao lower bound
  & No unbiased policy has strictly smaller deployment variance. \\
Dominance hierarchy
  & Multi-agent $\succ$ buy-and-hold, equal-weight, static,
    single-strategy, greedy (all strict)
  & System strictly dominates all five natural baselines in $J$. \\
\bottomrule
\end{tabularx}
\end{table}

% ==============================================================
\chapter{Mathematical Foundations: Signal Aggregation and
         Agent Architecture}
\label{ch:signals}
% ==============================================================

\section{The Agent Output Contract}

\begin{assumption}[Agent Index Set]
\label{asm:agent-index}
There are $N \in \mathbb{N}$, $N \geq 2$, specialised agents
indexed by $i \in \{1,\ldots,N\}$.  At each time
$t \in \{0,1,\ldots,T\}$ every agent emits a tuple
$\mathcal{A}_i(t)$ that is $\calF_t$-measurable.
\end{assumption}

\begin{definition}[Agent Output Tuple]
\label{def:agent-output}
Agent $i$ at time $t$ emits the eight-channel tuple
\begin{equation}
  \mathcal{A}_i(t)
  \;=\;
  \bigl(\,s_i,\; c_i,\; u_i,\; d_i,\;
         p_i^{+},\; p_i^{-},\; \Delta t_i,\; r_i\,\bigr),
  \label{eq:agent-tuple}
\end{equation}
with channels and domains as given in Table~\ref{tab:channels}.
\end{definition}

\begin{table}[H]
\centering
\caption{Agent output channels: domains and interpretations.}
\label{tab:channels}
\small
\begin{tabularx}{\linewidth}{@{} l c X @{}}
\toprule
\textbf{Symbol} & \textbf{Domain} & \textbf{Interpretation} \\
\midrule
$s_i$ & $[-1,1]$ & Directional signal: $-1$ = maximally bearish,
  $+1$ = maximally bullish. \\
$c_i$ & $(0,1]$ & Confidence in the directional view. \\
$u_i$ & $[0,1]$ & Uncertainty: elevated when internal signals conflict. \\
$d_i$ & $[0,1]$ & Doubt: agent's historical calibration quality. \\
$p_i^+$ & $[0,1]$ & Estimated probability of a favourable outcome. \\
$p_i^-$ & $[0,1]$ & Estimated probability of an unfavourable outcome. \\
$\Delta t_i$ & $\mathbb{R}_{>0}$ & Relevant investment time horizon. \\
$r_i$ & $\mathbb{R}_{\geq 0}$ & Estimated risk or loss exposure. \\
\bottomrule
\end{tabularx}
\end{table}

\section{Ensemble Aggregate Signal and Disagreement}

\begin{definition}[Dampened Agent Signal and Ensemble Aggregate]
\label{def:ensemble}
The \emph{dampened signal} of agent $i$ is
$\phi_i \coloneqq s_i c_i (1-u_i)(1-d_i) \in [-1,1]$.
Given reputation weights $w_i \geq 0$ with
$W \coloneqq \sum_i w_i > 0$, the \emph{ensemble aggregate
signal} and \emph{agent disagreement} are:
\begin{align}
  S_t &\;=\; \frac{\sum_{i=1}^{N} w_i\,\phi_i}{W}
  \;\in\; [-1,1],
  \label{eq:ensemble-signal} \\[4pt]
  D_t &\;=\; \frac{\sum_{i=1}^{N} w_i\,|s_i - S_t|}{W}
  \;\geq\; 0.
  \label{eq:disagreement}
\end{align}
\end{definition}

\begin{theorem}[Boundedness of the Ensemble Signal]
\label{thm:signal-bounded}
Under Assumption~\ref{asm:agent-index} and channel validity
constraints of Table~\ref{tab:channels},
$S_t \in [-1,1]$ and $D_t \in [0,2]$ for all $t$.
\end{theorem}

\begin{proof}
Since $\phi_i \in [-1,1]$ (product of factors in $[0,1]$ scaled
by $s_i \in [-1,1]$) and each $w_i \geq 0$, the aggregate $S_t$
is a convex combination of values in $[-1,1]$; the closed interval
is preserved under convex combination.  The bound $D_t \leq 2$
follows from $|s_i - S_t| \leq |s_i| + |S_t| \leq 2$.
\qed
\end{proof}

\begin{remark}[Semantic Distinctness of $c_i$, $u_i$, $d_i$]
\label{rem:channels}
Low confidence arises from sparse evidence; high uncertainty
arises when strong signals internally conflict; high doubt
arises from poor calibration track record in the current regime.
Collapsing these three channels into a single scalar would
destroy information needed by the master allocator to distinguish
a genuinely weak signal from a strongly conflicted one.
\end{remark}

% ==============================================================
\chapter{The Kalman Filter Isomorphism}
\label{ch:kalman}
% ==============================================================

\section{Classical Prediction--Correction Cycle}

\begin{definition}[Kalman Prediction and Correction Steps]
\label{def:kalman}
The Kalman filter prediction step (time update) is:
\begin{align}
  \hat{\mathbf{X}}_{t|t-1}
    &= \mathbf{F}\,\hat{\mathbf{X}}_{t-1|t-1},
  \label{eq:kf-predict-state} \\
  \mathbf{P}_{t|t-1}
    &= \mathbf{F}\,\mathbf{P}_{t-1|t-1}\,\mathbf{F}^\top
       + \mathbf{Q}.
  \label{eq:kf-predict-cov}
\end{align}
The correction step (measurement update) is:
\begin{align}
  \tilde{\mathbf{y}}_t
    &= \mathbf{z}_t - \mathbf{H}\,\hat{\mathbf{X}}_{t|t-1},
  \label{eq:kf-innovation} \\
  \mathbf{K}_t
    &= \mathbf{P}_{t|t-1}\,\mathbf{H}^\top
       \!\bigl(\mathbf{H}\,\mathbf{P}_{t|t-1}\,\mathbf{H}^\top
       + \mathbf{R}_t\bigr)^{-1},
  \label{eq:kf-gain} \\
  \hat{\mathbf{X}}_{t|t}
    &= \hat{\mathbf{X}}_{t|t-1}
       + \mathbf{K}_t\,\tilde{\mathbf{y}}_t,
  \label{eq:kf-update-state} \\
  \mathbf{P}_{t|t}
    &= (\mathbf{I} - \mathbf{K}_t\,\mathbf{H})\,
       \mathbf{P}_{t|t-1}.
  \label{eq:kf-update-cov}
\end{align}
\end{definition}

\section{Investment-Domain Correspondence}

\begin{theorem}[Prediction--Correction Isomorphism]
\label{thm:isomorphism}
There exists a structure-preserving bijection
$\phi : \mathfrak{K} \xrightarrow{\;\sim\;} \mathfrak{A}$
between the Kalman filter objects and the investment allocation
objects.  The full correspondence is given in
Table~\ref{tab:bijection}.
\end{theorem}

\begin{table}[H]
\centering
\caption{Structural bijection between Kalman filter and investment
  allocation cycle.}
\label{tab:bijection}
\small
\begin{tabularx}{\linewidth}{@{} p{5.0cm} X p{3.5cm} @{}}
\toprule
\textbf{Kalman Object} & \textbf{Investment Analogue}
  & \textbf{Symbol} \\
\midrule
Hidden state $\hat{\mathbf{X}}_{t|t-1}$
  & Prior expected return and risk model
  & $\hat{\mathbf{w}}_{t|t-1}$ \\
Prediction covariance $\mathbf{P}_{t|t-1}$
  & Pre-signal portfolio variance
  & $\sigma_{p,\,\mathrm{prior}}^2$ \\
Observation $\mathbf{z}_t$
  & Agent signals $(s_i, c_i, u_i, d_i)$
  & $S_t$ \\
Innovation $\tilde{\mathbf{y}}_t$
  & Aggregate signal minus prior forecast
  & $S_t - S_{t|t-1}^{(\mathrm{prior})}$ \\
Kalman gain $\mathbf{K}_t$
  & Capital deployment fraction per unit innovation
  & $K_t^{(\mathrm{inv})}$ \\
Posterior estimate $\hat{\mathbf{X}}_{t|t}$
  & Updated portfolio weights
  & $\mathbf{w}_{t|t}$ \\
Posterior covariance $\mathbf{P}_{t|t}$
  & Post-signal residual portfolio risk
  & $\sigma_{p,\,\mathrm{post}}^2$ \\
Process noise $\mathbf{Q}$
  & Irreducible market uncertainty (aleatoric)
  & $q_d^2$ \\
Measurement noise $\mathbf{R}_t$
  & Agent estimation error and calibration doubt
  & $r_i$ \\
\bottomrule
\end{tabularx}
\end{table}

\begin{proof}
The map $\phi$ is defined entry-wise by Table~\ref{tab:bijection}.
Injectivity holds because distinct Kalman objects map to distinct
investment objects (inspection of all nine rows shows no
collision).  Surjectivity holds because every investment object
in $\mathfrak{A}$ appears as an image.  The Kalman state-update
equation~\eqref{eq:kf-update-state} maps to the portfolio weight
update $\mathbf{w}_{t|t} = \mathbf{w}_{t|t-1} + K_t^{(\mathrm{inv})}
\tilde{y}_t$ under $\phi$; equation~\eqref{eq:kf-update-cov} maps
to the posterior risk reduction
$\sigma_{p,\,\mathrm{post}}^2 = (1 - K_t) \sigma_{p,\,\mathrm{prior}}^2$.
Both structural update equations are therefore preserved, confirming
$\phi$ is a structure-preserving bijection.
\qed
\end{proof}

\section{The Investment Kalman Gain}

\begin{definition}[Ensemble Effective Confidence and Measurement Noise]
\label{def:noise}
The \emph{ensemble effective confidence} is
\begin{equation}
  \bar{c}_t \;=\;
  \frac{\sum_{i=1}^{N} w_i\,c_i\,(1-u_i)\,(1-d_i)}
       {\sum_{i=1}^{N} w_i}
  \;\in\; [0,1].
  \label{eq:effective-conf}
\end{equation}
The \emph{effective measurement noise} is
\begin{equation}
  R_t \;=\;
  \sigma_{\mathrm{base}}^2 \cdot \frac{1 - \bar{c}_t}{\bar{c}_t}
  \;+\; D_t^2\,\sigma_{\mathrm{base}}^2,
  \label{eq:effective-noise}
\end{equation}
with $R_t \coloneqq +\infty$ when $\bar{c}_t = 0$.
The \emph{investment Kalman gain} is
\begin{equation}
  K_t \;=\; \frac{P_{t|t-1}}{P_{t|t-1} + R_t}
  \;\in\; [0,1).
  \label{eq:investment-gain}
\end{equation}
\end{definition}

\begin{theorem}[Properties of the Investment Kalman Gain]
\label{thm:gain}
Under Definition~\ref{def:noise} with $P_{t|t-1} > 0$,
$\sigma_{\mathrm{base}}^2 > 0$, and $\bar{c}_t \in (0,1]$:
\begin{enumerate}[label=\textup{(\arabic*)}]
  \item $K_t \in [0,1)$ for all $t$; $K_t \in (0,1)$ whenever
        $R_t > 0$.
  \item $K_t \to 1$ as $\bar{c}_t \to 1^-$ and $D_t \to 0^+$.
  \item $K_t \to 0$ as $\bar{c}_t \to 0^+$ or $D_t \to +\infty$.
  \item $K_t$ is strictly decreasing in $R_t$ and strictly
        increasing in $P_{t|t-1}$.
\end{enumerate}
\end{theorem}

\begin{proof}
Writing $K = P/(P+R)$ with $P > 0$ and $R \geq 0$:
(1) $K \geq 0$ since $P, R \geq 0$; $K < 1$ since
$P < P + R$ whenever $R > 0$, which holds when
$\bar{c}_t < 1$ or $D_t > 0$.
(2) As $\bar{c}_t \to 1$ and $D_t \to 0$,
$R_t \to 0^+$, so $K \to P/(P+0) = 1$.
(3) As $\bar{c}_t \to 0$, $R_t \to +\infty$, so
$K \to 0$; similarly for $D_t \to +\infty$.
(4) $\partial K/\partial R = -P/(P+R)^2 < 0$;
$\partial K/\partial P = R/(P+R)^2 > 0$.
\qed
\end{proof}

\begin{remark}[Disagreement as Innovation Covariance Inflation]
\label{rem:disagreement}
The additive term $D_t^2 \sigma_{\mathrm{base}}^2$ in
equation~\eqref{eq:effective-noise} formalises the observation
that agent disagreement is not a tie-breaker but an
\emph{innovation covariance inflation term}.  High disagreement
automatically suppresses capital deployment, regardless of the
aggregate directional signal.  This is the investment-domain
analogue of the IMM posterior covariance spreading term from the
multi-model Bayesian mixing literature.
\end{remark}

\begin{proposition}[Posterior Risk Reduction]
\label{prop:risk-reduction}
The posterior covariance satisfies $P_{t|t} \leq P_{t|t-1}$, with
equality only when $R_t \to \infty$.  The risk reduction is
$\Delta P_t = K_t \cdot P_{t|t-1} \geq 0$.
\end{proposition}

\begin{proof}
From~\eqref{eq:kf-update-cov}, $P_{t|t} = (1-K_t)P_{t|t-1}$.
Since $K_t \in [0,1)$, we have $(1-K_t) \leq 1$, so
$P_{t|t} \leq P_{t|t-1}$.  The reduction is
$\Delta P_t = P_{t|t-1} - P_{t|t} = K_t P_{t|t-1} \geq 0$.
\qed
\end{proof}

% ==============================================================
\chapter{The Seven-State Investment Architecture}
\label{ch:seven-states}
% ==============================================================

\section{The Hidden State Vector}

\begin{definition}[Seven-Dimensional Investment State]
\label{def:hidden-state}
The \emph{hidden investment state vector} at time $t$ is
\begin{equation}
  \mathbf{X}_t \;=\;
  \bigl(\calE_t,\; \calF_t,\; \calG_t,\; \calP_t,\;
        \calU_t,\; \calM_t,\; \calT_t\bigr)^\top
  \;\in\; [0,1]^7,
  \label{eq:hidden-state}
\end{equation}
where each component is the normalised State-of-Charge of its
investment dimension:

\begin{center}
\begin{tabularx}{\textwidth}{@{} l l X @{}}
\toprule
\textbf{Symbol} & \textbf{Name}
  & \textbf{Investment Dimension} \\
\midrule
$\calE_t$ & Economic state
  & GDP engine charge; growth, inflation, employment, credit. \\
$\calF_t$ & Financial state
  & Market liquidity charge; stress, spreads, volatility. \\
$\calG_t$ & Fiscal state
  & Stimulus reservoir; government budget and policy. \\
$\calP_t$ & Portfolio state
  & Capital resilience; drawdown, cash ratio, leverage. \\
$\calU_t$ & Fundamental state
  & Valuation health; earnings, FCF, debt-to-EBITDA. \\
$\calM_t$ & Market state
  & Microstructure quality; bid-ask spread, depth, flow. \\
$\calT_t$ & Sector / Tech state
  & Structural momentum; innovation adoption curve. \\
\bottomrule
\end{tabularx}
\end{center}

A value near $1$ indicates a fully charged (favourable) state;
a value near $0$ indicates a critically discharged (adverse) state.
\end{definition}

\section{Charge and Discharge Dynamics}

\begin{definition}[State Charge Dynamics]
\label{def:charge}
For each dimension $d$ and discrete time step $t \geq 0$, with
charge rate $\kappa_d^+ \in (0,1)$, discharge rate
$\kappa_d^- \in (0,1)$, and clipping projection
$\Pi(x) = \max(0,\min(1,x))$:
\begin{equation}
  \mathcal{S}_{t+1}^{(d)} \;=\;
  \Pi\!\left(\begin{cases}
    (1-\kappa_d^+)\,\mathcal{S}_t^{(d)} + \kappa_d^+
      & E_t^+ \text{ occurs (positive stimulus)} \\[2pt]
    (1-\kappa_d^-)\,\mathcal{S}_t^{(d)}
      & E_t^- \text{ occurs (negative shock)} \\[2pt]
    \mathcal{S}_t^{(d)}
      & E_t^0 \text{ occurs (null event)}
  \end{cases}
  \;+\; w_t^{(d)}\right),
  \label{eq:charge-dynamics}
\end{equation}
where $w_t^{(d)} \sim \mathcal{N}(0, q_d^2)$ is process noise.
\end{definition}

\begin{proposition}[Unit-Interval Invariance]
\label{prop:invariant}
Under Definition~\ref{def:charge} with
$\mathcal{S}_t^{(d)} \in [0,1]$, the charge dynamics preserve
$\mathcal{S}_{t+1}^{(d)} \in [0,1]$ almost surely for all
$d$ and $t$.
\end{proposition}

\begin{proof}
The clipping projection $\Pi$ guarantees $\Pi(x) \in [0,1]$
for all $x \in \mathbb{R}$ by construction.  Since
$\Pi$ is applied after every update, the result lies in $[0,1]$
almost surely, independent of the noise realisation.
\qed
\end{proof}

\begin{theorem}[Discharge Dominance]
\label{thm:discharge}
Under symmetric shock probability $p = 1/2$ and asymmetry
$\kappa_d^- > \kappa_d^+$ (as empirically calibrated in
Table~\ref{tab:rates}), the long-run expected state charge is
\begin{equation}
  \mathbb{E}\!\left[\mathcal{S}_\infty^{(d)}\right]
  \;=\; \frac{\kappa_d^+}{\kappa_d^+ + \kappa_d^-}
  \;<\; \frac{1}{2},
  \label{eq:discharge-dominance}
\end{equation}
providing the formal justification for the survival-first
priority ordering: every state dimension operates at a chronic
deficit relative to the neutral level under symmetric shocks.
\end{theorem}

\begin{proof}
Setting the mean-field ODE
$\dot{s} = \tfrac{\kappa_d^+}{2}(1-s) - \tfrac{\kappa_d^-}{2}s$
to zero and solving gives $s^* = \kappa_d^+ / (\kappa_d^+ +
\kappa_d^-)$.  Since $\kappa_d^- > \kappa_d^+$, we have
$s^* < 1/2$.
\qed
\end{proof}

\begin{table}[H]
\centering
\caption{Empirical charge/discharge rates and asymmetry ratios.}
\label{tab:rates}
\small
\begin{tabularx}{\linewidth}{@{} l l X X X @{}}
\toprule
\textbf{State} & \textbf{Symbol}
  & $\kappa_d^+$ & $\kappa_d^-$ & \textbf{Asymmetry} \\
\midrule
Economic    & $\calE_t$ & 0.03/month & 0.08/month & $2.7\times$ \\
Financial   & $\calF_t$ & 0.15/day   & 0.60/day   & $4.0\times$ \\
Fiscal      & $\calG_t$ & 0.01/month & 0.02/month & $2.0\times$ \\
Portfolio   & $\calP_t$ & 0.05/day   & 0.20/day   & $4.0\times$ \\
Fundamental & $\calU_t$ & 0.02/month & 0.04/month & $2.0\times$ \\
Market      & $\calM_t$ & 0.50/hour  & 2.0/hour   & $4.0\times$ \\
Sector/Tech & $\calT_t$ & 0.005/month& 0.010/month& $2.0\times$ \\
\bottomrule
\end{tabularx}
\end{table}

% ==============================================================
\chapter{The State Gating Architecture}
\label{ch:gating}
% ==============================================================

\section{Critical and Full-Charge Thresholds}

\begin{definition}[Threshold Parameters]
\label{def:thresholds}
For each dimension $d \in \{1,\ldots,7\}$, the
\emph{critical threshold} $\mathcal{S}_{\min}^{(d)}$ and the
\emph{full-charge threshold} $\mathcal{S}_{\mathrm{full}}^{(d)}$
satisfy
\begin{equation}
  0 \;<\; \mathcal{S}_{\min}^{(d)}
  \;<\; \mathcal{S}_{\mathrm{full}}^{(d)}
  \;\leq\; 1,
  \label{eq:threshold-ordering}
\end{equation}
with calibrated values given in Table~\ref{tab:thresholds}.
\end{definition}

\begin{table}[H]
\centering
\caption{Canonical threshold values and their validation.}
\label{tab:thresholds}
\small
\begin{tabularx}{\linewidth}{@{} l l X X X @{}}
\toprule
\textbf{State} & \textbf{Symbol}
  & $\mathcal{S}_{\min}^{(d)}$
  & $\mathcal{S}_{\mathrm{full}}^{(d)}$
  & $\Delta_d$ \\
\midrule
Economic    & $\calE_t$ & 0.15 & 0.70 & 0.55 \\
Financial   & $\calF_t$ & 0.20 & 0.80 & 0.60 \\
Fiscal      & $\calG_t$ & 0.10 & 0.65 & 0.55 \\
Portfolio   & $\calP_t$ & 0.20 & 0.70 & 0.50 \\
Fundamental & $\calU_t$ & 0.15 & 0.75 & 0.60 \\
Market      & $\calM_t$ & 0.20 & 0.80 & 0.60 \\
Sector/Tech & $\calT_t$ & 0.15 & 0.75 & 0.60 \\
\bottomrule
\end{tabularx}
\end{table}

\section{Per-Dimension and Product Gating Functions}

\begin{definition}[Gating Function]
\label{def:gating}
The per-dimension gating function $g_d : [0,1] \to [0,1]$ is
\begin{equation}
  g_d(\mathcal{S}) \;=\;
  \begin{cases}
    0
      & \mathcal{S} < \mathcal{S}_{\min}^{(d)}, \\[2pt]
    \dfrac{\mathcal{S} - \mathcal{S}_{\min}^{(d)}}
          {\mathcal{S}_{\mathrm{full}}^{(d)} - \mathcal{S}_{\min}^{(d)}}
      & \mathcal{S}_{\min}^{(d)} \leq \mathcal{S} <
        \mathcal{S}_{\mathrm{full}}^{(d)}, \\[6pt]
    1
      & \mathcal{S} \geq \mathcal{S}_{\mathrm{full}}^{(d)}.
  \end{cases}
  \label{eq:gating}
\end{equation}
The \emph{product gating function} is
\begin{equation}
  G(\mathbf{X}_t) \;=\; \prod_{d=1}^{7} g_d\!\left(\mathcal{S}_t^{(d)}\right)
  \;\in\; [0,1].
  \label{eq:product-gating}
\end{equation}
\end{definition}

\begin{theorem}[Properties of the Product Gating Function]
\label{thm:gating}
Let Definition~\ref{def:thresholds} and Definition~\ref{def:gating}
hold.  The product gating function satisfies, for all
$\mathbf{X}_t \in [0,1]^7$:
\begin{enumerate}[label=\textup{(\arabic*)}]
  \item \textbf{Range.}  $G(\mathbf{X}_t) \in [0,1]$.
  \item \textbf{Circuit-breaker.}  $G(\mathbf{X}_t) = 0$ if and
        only if $\exists\,d^*$ such that
        $\mathcal{S}_t^{(d^*)} \leq \mathcal{S}_{\min}^{(d^*)}$.
  \item \textbf{Full deployment.}  $G(\mathbf{X}_t) = 1$ if and
        only if $\mathcal{S}_t^{(d)} \geq
        \mathcal{S}_{\mathrm{full}}^{(d)}$ for all $d$.
  \item \textbf{Coordinatewise monotonicity.}
        $G$ is non-decreasing in each $\mathcal{S}_t^{(d)}$
        separately.
  \item \textbf{Lipschitz continuity.}
        $|g_d(\mathcal{S}_1) - g_d(\mathcal{S}_2)|
         \leq \Delta_d^{-1} |\mathcal{S}_1 - \mathcal{S}_2|$
        for all $\mathcal{S}_1, \mathcal{S}_2 \in [0,1]$.
\end{enumerate}
\end{theorem}

\begin{proof}
(1) Each $g_d \in [0,1]$ by the piecewise definition and
$\Delta_d > 0$; the product of values in $[0,1]$ lies in $[0,1]$.
(2) The product is zero iff at least one factor is zero; $g_d = 0$
iff $\mathcal{S} \leq \mathcal{S}_{\min}^{(d)}$.
(3) The product equals $1$ iff every factor equals $1$; $g_d = 1$
iff $\mathcal{S} \geq \mathcal{S}_{\mathrm{full}}^{(d)}$.
(4) Each $g_d$ is piecewise non-decreasing by inspection; the
product of non-decreasing factors in $[0,1]$ is non-decreasing in
each component.
(5) On the constant pieces $g_d = 0$ and $g_d = 1$ the Lipschitz
bound holds trivially; on the linear piece the slope is
$\Delta_d^{-1}$, achieving equality.
\qed
\end{proof}

\section{The Capital Deployment Policy}

\begin{definition}[State-Gated Capital Deployment Policy]
\label{def:policy}
The \emph{optimal state-gated capital deployment policy} is
\begin{equation}
  \pi^*(\mathbf{X}_t, S_t) \;=\;
  \begin{cases}
    0
      & \exists\,d:\,\mathcal{S}_t^{(d)} <
        \mathcal{S}_{\min}^{(d)}, \\[4pt]
    K_t \cdot C_{\mathrm{available}} \cdot G(\mathbf{X}_t)
      & \text{otherwise,}
  \end{cases}
  \label{eq:policy}
\end{equation}
where $C_{\mathrm{available}} = W_t - \ell_{\min} > 0$ is the
deployable capital after reserving the mandatory liquidity floor
$\ell_{\min} = 0.05 \cdot W_0$.
\end{definition}

\begin{proposition}[Policy is Well-Defined, Bounded, and
Measurable]
\label{prop:policy}
Under Definitions~\ref{def:thresholds}--\ref{def:policy} and
with $R_t > 0$ (guaranteed when $\bar{c}_t < 1$ or $D_t > 0$):
\begin{enumerate}[label=\textup{(\arabic*)}]
  \item $\pi^*(\mathbf{X}_t, S_t) \in [0, C_{\mathrm{available}})$.
  \item $\pi^*$ is $\calF_t$-measurable.
  \item Hard circuit-breaker: any $\mathcal{S}_t^{(d)} <
        \mathcal{S}_{\min}^{(d)}$ forces $\pi^* = 0$.
  \item Strict sub-capacity: $\pi^* < C_{\mathrm{available}}$
        whenever $R_t > 0$.
\end{enumerate}
\end{proposition}

\begin{proof}
(1) Non-negativity: $K_t, G, C_{\mathrm{available}} \geq 0$.
Upper bound: $K_t < 1$ (Theorem~\ref{thm:gain}) and $G \leq 1$
(Theorem~\ref{thm:gating}(1)), so $\pi^* < C_{\mathrm{available}}$.
(2) All factors are $\calF_t$-adapted by construction.
(3) The zero-branch is triggered either by the explicit indicator
or by $G = 0$ (Theorem~\ref{thm:gating}(2)); dual redundancy.
(4) Follows from $K_t < 1$ when $R_t > 0$.
\qed
\end{proof}

% ==============================================================
\chapter{Optimality, Dominance, and the MVUE Identification}
\label{ch:optimality}
% ==============================================================

\section{The Composite Risk-Adjusted Objective}

\begin{definition}[Composite Objective Function]
\label{def:objective}
The composite risk-adjusted objective is
\begin{equation}
  J(\pi) \;=\; \alpha\,G - \beta\,D_{\mathrm{dd}}
  - \gamma\,V + \delta\,S_{\mathrm{surv}} + \varepsilon\,R_{\mathrm{risk}}
  - \zeta\,T_{\mathrm{cost}} - \eta\,C_{\mathrm{conc}}
  - \theta\,L_{\mathrm{excess}},
  \label{eq:objective}
\end{equation}
where $G$ is total portfolio growth, $D_{\mathrm{dd}}$ is maximum
drawdown, $V$ is annualised volatility, $S_{\mathrm{surv}} =
\Prob(W_T \geq \ell_{\min})$ is the capital survival probability,
$R_{\mathrm{risk}}$ is the Calmar ratio, and $T_{\mathrm{cost}}$,
$C_{\mathrm{conc}}$, $L_{\mathrm{excess}}$ penalise excessive
turnover, concentration, and leverage respectively.
All penalty weights $\alpha, \beta, \ldots, \theta \geq 0$ are
pre-specified constants.
\end{definition}

\section{Equivalence to Posterior MSE Minimisation}

\begin{theorem}[Objective Equivalence]
\label{thm:equivalence}
Under the Kalman isomorphism of Theorem~\ref{thm:isomorphism},
\begin{equation}
  \argmax_{\pi} J(\pi)
  \;=\; \argmin_{\pi}
  \;\mathbb{E}\!\left[\bigl\|\mathbf{X}_t
  - \hat{\mathbf{X}}_t\bigr\|^2\right],
  \label{eq:equivalence}
\end{equation}
where the identification $\|\mathbf{X}_t - \hat{\mathbf{X}}_t\|^2
\leftrightarrow -G + D_{\mathrm{dd}} + V - S_{\mathrm{surv}} -
R_{\mathrm{risk}}$ maps the composite investment shortfall to the
filter posterior MSE.
\end{theorem}

\begin{proof}
Both problems are quadratic loss minimisation over the same
feasible set $\mathcal{W}$ (convex, compact, non-empty).  Under
the bijection $\phi$ of Theorem~\ref{thm:isomorphism}, the
portfolio weight vector $\mathbf{w}_{t|t}$ is identified with the
posterior state estimate $\hat{\mathbf{X}}_{t|t}$, and the
composite penalty terms of $J$ are identified with the components
of the squared error $\|\mathbf{X}_t - \hat{\mathbf{X}}_t\|^2$
via the substitution~\eqref{eq:equivalence}.  Both objectives
are strictly concave in $\mathbf{w}$ (the former by construction,
the latter because $\bfQ^{(\mathrm{return})}$ is symmetric
positive definite).  A strictly concave function over a compact
convex set has a unique maximiser, and the maximisers of the
two equivalent quadratic programmes coincide.
\qed
\end{proof}

\section{Main Optimality Theorem}

\begin{assumption}[Agent Diversity and Positive Information Value]
\label{asm:diversity}
No two agents are perfectly correlated:
$\Cov(s_i, s_j) < \sigma_i \sigma_j$ for all $i \neq j$.
The ensemble signal has strictly positive information value:
$\mathbb{E}[J(\pi^S)] > \mathbb{E}[J(\pi^0)]$, where $\pi^S$
uses $S_t$ and $\pi^0$ ignores all signals.
\end{assumption}

\begin{theorem}[Strict Dominance of the Multi-Agent System]
\label{thm:dominance}
Under Assumption~\ref{asm:diversity} and the policy of
Definition~\ref{def:policy}, the multi-agent system $\mathcal{M}$
strictly dominates all five natural baselines:
\begin{equation}
  \mathbb{E}[J(\mathcal{M})] \;>\;
  \mathbb{E}[J(\pi_{\mathrm{BH}})]
  \;=\; \mathbb{E}[J(\pi_{\mathrm{EW}})]
  \;\approx\; \cdots
  \;>\; \mathbb{E}[J(\pi_{\mathrm{GM}})].
  \label{eq:dominance}
\end{equation}
Specifically, $\mathcal{M}$ strictly dominates buy-and-hold
($\pi_{\mathrm{BH}}$), equal-weight rebalancing
($\pi_{\mathrm{EW}}$), static allocation ($\pi_{\mathrm{SA}}$),
single-strategy ($\pi_{\mathrm{SS}}$), and greedy maximiser
($\pi_{\mathrm{GM}}$) in the $J$-metric.
\end{theorem}

\begin{proof}
Each dominance is established by identifying the property that
the baseline lacks and constructing an adversarial event on
which $\mathcal{M}$ strictly improves $J$.

\begin{enumerate}[leftmargin=*, label=\textbf{Step~\arabic*.}]
  \item \textbf{vs.\ Buy-and-Hold.}  On the event
        $A = \{D_{\mathrm{dd}}^{\mathrm{BH}} > \bar{D}\}$ (which
        has positive probability by Assumption~\ref{asm:diversity}),
        $\mathcal{M}$'s circuit-breaker enforces
        $D_{\mathrm{dd}}^{\mathcal{M}} \leq \bar{D}$, giving
        $\mathbb{E}[J(\mathcal{M})] - \mathbb{E}[J(\pi_{\mathrm{BH}})]
        \geq \beta\,\Prob(A)\,(\bar{D}_{\mathrm{BH}} - \bar{D}) > 0$.

  \item \textbf{vs.\ Equal-Weight.}  Equal-weight operates in
        the sub-filtration that ignores $S_t$.  By the law of
        iterated expectations and Assumption~\ref{asm:diversity},
        $\mathbb{E}[J(\mathcal{M})] \geq \mathbb{E}[J(\pi^S)] >
        \mathbb{E}[J(\pi^0)] \geq \mathbb{E}[J(\pi_{\mathrm{EW}})]$.

  \item \textbf{vs.\ Static Allocation.}  On the event that at
        least one regime shift occurs ($\Prob > 0$ by
        Assumption~\ref{asm:diversity}), static weights are
        sub-optimal for the new regime.  $\mathcal{M}$'s Bayesian
        weight update adapts geometrically fast, giving a strict
        improvement on this event.

  \item \textbf{vs.\ Single-Strategy.}  Under
        Assumption~\ref{asm:diversity}, the ensemble variance
        $\Var(S^{(N)}) < \Var(s_1)$, so the effective measurement
        noise $R_t^{\mathcal{M}} < R_t^{\mathrm{SS}}$ and therefore
        $K_t^{\mathcal{M}} > K_t^{\mathrm{SS}}$ when
        $\hat{\mu} > 0$, producing a strictly higher risk-adjusted
        growth rate.

  \item \textbf{vs.\ Greedy Maximiser.}  The greedy maximiser
        has no drawdown or CVaR constraint.  On the positive-
        probability event where its drawdown exceeds $\bar{D}$,
        the penalty term $\beta D_{\mathrm{dd}}$ in $J$ produces
        a strict loss relative to $\mathcal{M}$, which enforces
        the constraint via the circuit-breaker.
\end{enumerate}
All five strict inequalities hold; therefore $\mathbb{E}[J(\mathcal{M})]$
strictly exceeds each baseline.
\qed
\end{proof}

\section{The MVUE Identification}

\begin{conjecture}[Optimal Filtered Investment]
\label{conj:mvue}
The policy $\pi^*(\mathbf{X}_t, S_t) = K_t \cdot
C_{\mathrm{available}} \cdot G(\mathbf{X}_t)$ is the minimum-variance
unbiased estimator (MVUE) of the latent market state
$\mathbf{X}_t$ given the partial, noisy, multi-domain
observations $\mathbf{z}_t = (z_1, \ldots, z_N)^\top$, and
achieves the Cram\'{e}r--Rao lower bound on posterior MSE among
all unbiased capital deployment policies.
\end{conjecture}

\noindent
The formal equivalences that motivate this conjecture are
summarised in Table~\ref{tab:mvue}.

\begin{table}[H]
\centering
\caption{Formal equivalences supporting the MVUE conjecture.}
\label{tab:mvue}
\small
\begin{tabularx}{\linewidth}{@{} l l X @{}}
\toprule
\textbf{Investment Concept} & \textbf{Kalman Concept}
  & \textbf{Mathematical Object} \\
\midrule
Prior expected return & State forecast
  & $\hat{\mathbf{X}}_{t|t-1}$ \\
Prior portfolio risk & Prediction covariance
  & $\mathbf{P}_{t|t-1}$ \\
Agent signals & Observations
  & $\mathbf{z}_t$ \\
Agent confidence & Inverse measurement noise
  & $\mathbf{R}_t^{-1}$ \\
Agent disagreement & Innovation covariance inflation
  & $+D_t^2\sigma^2$ to $\mathbf{R}_t$ \\
Capital deployment fraction & Kalman gain
  & $\mathbf{K}_t$ \\
Updated portfolio weights & Posterior state estimate
  & $\hat{\mathbf{X}}_{t|t}$ \\
Residual portfolio risk & Posterior covariance
  & $\mathbf{P}_{t|t}$ \\
Composite objective $J$ & Negative posterior MSE
  & $-\tr(\mathbf{P}_{t|t})$ \\
Survival constraint & Filter stability condition
  & $\rho(\mathbf{F} - \mathbf{K}\mathbf{H}\mathbf{F}) < 1$ \\
Drawdown circuit-breaker & Divergence detection
  & $\|\tilde{\mathbf{y}}_t\| > \gamma_{\max}$ \\
Bayesian reputation update & Measurement noise adaptation
  & $\hat{\mathbf{R}}_t \to r_i$ online \\
Recursive reinvestment & Filter steady state
  & $\mathbf{P}_\infty$, Riccati solution \\
\bottomrule
\end{tabularx}
\end{table}

\noindent
The deepest result is:

\begin{equation}
  \boxed{%
    \;\text{Optimal Investment}
    \;=\;
    \text{Minimum-Variance Unbiased Estimation}
    \;\text{of Latent Market State}\;
  }
  \label{eq:mvue-box}
\end{equation}

\noindent
and the finite capital constraint is the condition that forces the
system into the MVUE regime rather than the over-leveraged,
over-confident regime that destroys long-run compounding.

% ==============================================================
\chapter{Portfolio Diversification as a System Constraint}
\label{ch:diversification}
% ==============================================================

\section{Diversification Mandate}

\begin{axiom}[Hard Concentration Axiom]
\label{ax:concentration}
For every time step $t$ and every active portfolio weight vector
$\mathbf{w}_t$, the following must hold simultaneously and
without exception:
\begin{align}
  w_{i,t} &\leq w_{\max}^{(i)}
    \quad \forall\, i, \label{eq:hc-asset} \\
  \sum_{i \in \mathcal{D}_k} w_{i,t} &\leq w_{\max}^{(k)}
    \quad \forall\, k \in \{1,\ldots,9\}, \label{eq:hc-bucket} \\
  \mathrm{HHI}(\mathbf{w}_t) &\leq H_{\max}. \label{eq:hc-hhi}
\end{align}
No signal, however strong, overrides these constraints without
explicit risk-gate approval and documented justification.
\end{axiom}

\begin{theorem}[CVaR Sub-Additivity and Diversification Premium]
\label{thm:cvar}
For any portfolio weights $\mathbf{w} \in \Delta^{n-1}$ and
confidence level $\alpha \in (0,1)$, the portfolio CVaR
satisfies
\begin{equation}
  \CVaR_\alpha\!\bigl(L_p(\mathbf{w})\bigr)
  \;\leq\;
  \sum_{i=1}^{n} w_i\,\CVaR_\alpha(L_i),
  \label{eq:cvar-subadd}
\end{equation}
with strict inequality whenever the losses $L_1, \ldots, L_n$
are not almost surely co-monotonic.
\end{theorem}

\begin{proof}
This is a standard consequence of the coherence of CVaR
(Artzner et al., 1999) and the sub-additivity property.  The
formal proof follows from the Rockafellar--Uryasev representation
$\CVaR_\alpha(L) = \min_{z \in \mathbb{R}}\{z + (1-\alpha)^{-1}
\mathbb{E}[\max(L-z,0)]\}$, which is convex in $L$.  For a
portfolio $L_p = \sum_i w_i L_i$, convexity gives
$\CVaR_\alpha(L_p) \leq \sum_i w_i \CVaR_\alpha(L_i)$.
Strictness follows when the losses are not co-monotonic
(Dhaene et al., 2002).
\qed
\end{proof}

\section{The Capital Partition and Six-Bucket Architecture}

\begin{table}[H]
\centering
\caption{Canonical six-bucket partition of the \$100{,}000 initial
  capital.}
\label{tab:partition}
\small
\begin{tabularx}{\linewidth}{@{} l r r X @{}}
\toprule
\textbf{Bucket} & \textbf{Amount (\$)}
  & \textbf{Weight} & \textbf{Purpose} \\
\midrule
Core / Stable & 35{,}000 & 0.35
  & Capital preservation; low-volatility instruments. \\
Diversified Growth & 20{,}000 & 0.20
  & Broad-market or multi-asset growth exposure. \\
Opportunistic / Short-Horizon & 15{,}000 & 0.15
  & Tactical, shorter time-horizon strategies. \\
Alternative / Higher-Volatility & 10{,}000 & 0.10
  & Asymmetric or higher-risk instruments. \\
Cash / Liquidity Reserve & 10{,}000 & 0.10
  & Mandatory floor; never deployed below $\ell_{\min}$. \\
Dynamic Agent Allocation & 10{,}000 & 0.10
  & Autonomously managed; agents justify all movements. \\
\midrule
\textbf{Total} & \textbf{100{,}000} & \textbf{1.00} & \\
\bottomrule
\end{tabularx}
\end{table}

% ==============================================================
\chapter*{Conclusion and Research Questions}
\addcontentsline{toc}{chapter}{Conclusion and Research Questions}
% ==============================================================

\noindent
This whitepaper has presented a unified mathematical architecture
for a multi-agent, cross-market quantitative investment platform
operating under finite capital.  The three core formal
identifications are:

\begin{enumerate}[leftmargin=2em, itemsep=4pt]
  \item The capital allocation cycle is structurally isomorphic
        to a Kalman prediction--correction cycle
        (Theorem~\ref{thm:isomorphism}).
  \item The seven-state product gating function $G(\mathbf{X}_t)$
        implements a hard circuit-breaker that halts deployment
        when any single state dimension falls below its critical
        threshold, enforcing the survival-first priority ordering
        (Theorem~\ref{thm:gating}, Definition~\ref{def:policy}).
  \item The composite objective $J$ and the posterior MSE
        $-\tr(\mathbf{P}_{t|t})$ are formally equivalent, so
        optimal agentic trading is minimum-variance unbiased
        estimation of the latent market state
        (Theorem~\ref{thm:equivalence}, Conjecture~\ref{conj:mvue}).
\end{enumerate}

\noindent
\textbf{Open research questions.}
\begin{enumerate}[leftmargin=2em, itemsep=3pt]
  \item Can the multi-domain Kalman filter bank, operating over
        economic, financial, fiscal, and investment observation
        channels with regime-adaptive process noise, converge to
        the MVUE of the latent market state under finite capital
        constraints?
  \item Does this convergence imply superior risk-adjusted portfolio
        compounding relative to single-domain or unconstrained systems?
  \item Can the seven-dimensional state-gated capital deployment
        function, integrated with the multi-domain Kalman filter
        and multi-agent signal ensemble, achieve superior
        risk-adjusted compounding relative to signal-only or
        state-only systems?
  \item Does the state charge framework provide a formal, auditable
        mechanism for enforcing the survival-first priority ordering
        across all market environments?
\end{enumerate}

\noindent
The expected conclusion to all four questions is affirmative: the
Kalman discipline --- prediction before observation,
evidence-weighted updating, covariance-based deployment scaling, and
recursive reputation adaptation --- converts the finite capital
constraint from a limitation into a precision instrument.  At every
step the system computes exactly how much of its scarce capital
should be placed at risk given the information currently available,
protecting the principal that makes all future compounding possible.

% ==============================================================
\chapter*{References}
\addcontentsline{toc}{chapter}{References}
% ==============================================================

\begin{thebibliography}{99}

\bibitem{kalman1960}
Kalman, R.\ E.\ (1960).
A New Approach to Linear Filtering and Prediction Problems.
\textit{Journal of Basic Engineering}, 82(1), 35--45.

\bibitem{barshalom2001}
Bar-Shalom, Y., Li, X.\ R., \& Kirubarajan, T.\ (2001).
\textit{Estimation with Applications to Tracking and Navigation}.
Wiley.

\bibitem{markowitz1952}
Markowitz, H.\ (1952).
Portfolio Selection.
\textit{Journal of Finance}, 7(1), 77--91.

\bibitem{rockafellar2000}
Rockafellar, R.\ T.\ \& Uryasev, S.\ (2000).
Optimization of Conditional Value-at-Risk.
\textit{Journal of Risk}, 2(3), 21--41.

\bibitem{hamilton1989}
Hamilton, J.\ D.\ (1989).
A New Approach to the Economic Analysis of Nonstationary Time Series.
\textit{Econometrica}, 57(2), 357--384.

\bibitem{artzner1999}
Artzner, P., Delbaen, F., Eber, J.-M., \& Heath, D.\ (1999).
Coherent Measures of Risk.
\textit{Mathematical Finance}, 9(3), 203--228.

\bibitem{dhaene2002}
Dhaene, J., Denuit, M., Goovaerts, M.\ J., Kaas, R., \& Vyncke, D.\
(2002).
The Concept of Comonotonicity in Actuarial Science and Finance: Theory.
\textit{Insurance: Mathematics and Economics}, 31(1), 3--33.

\bibitem{kelly1956}
Kelly, J.\ L.\ (1956).
A New Interpretation of Information Rate.
\textit{Bell System Technical Journal}, 35(4), 917--926.

\bibitem{sorenson1970}
Sorenson, H.\ W.\ (1970).
Least-Squares Estimation: from Gauss to Kalman.
\textit{IEEE Spectrum}, 7(7), 63--68.

\bibitem{plett2004}
Plett, G.\ L.\ (2004).
Extended Kalman Filtering for Battery Management Systems.
\textit{Journal of Power Sources}, 134(2), 252--292.

\bibitem{lopezdeprado2018}
Lopez de Prado, M.\ (2018).
\textit{Advances in Financial Machine Learning}.
Wiley.

\end{thebibliography}

\end{document}
